\clearpage
\fancyhead[L]{\small HINDS STERN. PARTHENOPE. VICTORIA.}
\SourceLeaf{483}
\begin{center}
\rule{0.55\linewidth}{0.4pt}\\[2pt]
Astronomische Nachrichten. Band XXVIII. Nr. 650. S. 23. 2. November 1848.\\[2pt]
\rule{0.55\linewidth}{0.4pt}
\end{center}
\begin{flushright}Göttingen, 19. August 1848.\end{flushright}

Ich konnte den Meridiankreis erst Ende Juli wieder aufstellen und nutzte die Gelegenheit, \textsc{Hind}s Stern noch im Meridian zu beobachten. Meine Beobachtung ergab:

\begin{center}
1848 Juli 30.\quad Rektasz. $16^{\rm u}\ 51^{\rm m}\ 1^{\rm s}77$\\
Südl. Dekl. $12^\circ\ 39'\ 10''6$
\end{center}

Professor \textsc{Goldschmidt} hatte aus früheren Beobachtungen bereits den mittleren Ort bestimmt

\begin{center}
für 1848 Jan. 1.\quad Rektasz. $16^{\rm u}\ 50^{\rm m}\ 59^{\rm s}06$\\
Dekl. $-12^\circ\ 39'\ 17''3$
\end{center}

und hat jetzt alle Beobachtungen, die uns bekannt geworden sind, damit verglichen; \mbox{---\quad---\quad---}

